\documentclass[../main.tex]{subfiles}

\begin{document}
\section{Experimental Results (RQ1 \& RQ2)} \label{exp_results}

The results of FinBERT, the baseline methods and state-of-the-art on Financial PhraseBank dataset classification task can be seen on table \ref{tab:results}. We present the result on both the whole dataset and subset with 100\% annotator agreement. 

    \subfile{tables/results.tex}
    
For all of the measured metrics, FinBERT performs clearly the best among both the methods we implemented ourselves (LSTM and ULMFit) and the models reported by other papers (LPS \cite{Malo2014}, HSC \cite{Krishnamoorthy2018}, FinSSLX \cite{Maia20182}). LSTM classifier with no language model information performs the worst. In terms of accuracy, it is close to LPS and HSC, (even better than LPS for examples with full agreement), however it produces a low F1-score. That is due to it performing much better in neutral class. LSTM classifier with ELMo embeddings improves upon LSTM with static embeddings in all of the measured metrics. It still suffers from low average F1-score due to poor performance in less represented labels. But it's performance is comparable with LPS and HSC, besting them in accuracy. So contextualized word embeddings produce close performance to machine learning based methods for dataset of this size. 

ULMFit significantly improves on all of the metrics and it doesn't suffer from model performing much better in some classes than the others. It also handily beats the machine learning based models LPS and HSC. This shows the effectiveness of language model pre-training. AWD-LSTM is a very large model and it would be expected to suffer from over-fitting with this small of a dataset. But due to language model pre-training and effective training strategies, it is able to overcome small data problem. ULMFit also outperforms FinSSLX, which has a text simplification step as well as pre-training of word embeddings on a large financial corpus with sentiment labels. 

\begin{figure}[h] 
  \centering
  \includegraphics[width=\linewidth]{images/loss.png}
  \caption{Test loss different training set sizes}
  \label{fig:loss}
  \Description{The 1907 Franklin Model D roadster.}
\end{figure}

FinBERT outperforms ULMFit, and consequently all of the other methods in all metrics. In order to measure the performance of the models on different sizes of labeled training datasets, we ran LSTM classifiers, ULMFit and FinBERT on 5 different configurations. The result can be seen on figure \ref{fig:loss}, where the cross entropy losses on test set for each model are drawn. 100 training examples is too low for all of the models. However, once the training size becomes 250, ULMFit and FinBERT starts to successfully differentiate between labels, with an accuracy as high as 80\% for FinBERT. All of the methods consistently get better with more data, but ULMFit and FinBERT does better with 250 examples than LSTM classifiers do with the whole dataset. This shows the effectiveness of language model pre-training.

The results for FiQA sentiment dataset, are presented on table \ref{tab:regression}. Our model outperforms state-of-the-art models for both MSE and $R^2$. It should be noted that the test set these two papers \cite{Yang2018} \cite{Piao2018} use is the official FiQA Task 1 test set. Since we don't have access to that we report the results on 10-Fold cross validation. There is no indication on \cite{Maia2018} that the train and test sets they publish come from different distributions and our model can be interpreted to be at disadvantage since we need to set aside a subset of training set as test set, while state-of-the-art papers can use the complete training set.
    \subfile{tables/regression.tex}
    
\end{document}
